\section{The geodesic space and its integral classes}
\label{geo:section}

Throughout this section, $M$ has quaternionic model type, so that
\[
 n=4r,\qquad r\ge2,\qquad d=4r-1,\qquad k=3,
 \qquad H^*(M;\Z)=\Z[u]/(u^{r+1}),\quad |u|=4.
\]
We retain the normalization $D=\pi/2$ and $T=2D=\pi$.
Choose the sign of $u$ as permitted by Theorem~\ref{sig:coefficient},
so that the actual four-plane bundle $\mathcal V$ on cut incidence
$F$ satisfies
\begin{equation}\label{geo:incidence-input}
 p_1(\mathcal V)=2(u_1+u_2).
\end{equation}
Orient $M$ by $\int_Mu^r=1$.
Only the characteristic-class conclusion
\eqref{geo:incidence-input} and the bundle identification in
Proposition~\ref{inc:normal} are used from the incidence argument.

Let $U=SM$, let $\pi:U\to M$ be the unit tangent projection, and
write $z=\pi^*u$.  The prime period is $T$, so the action
\[
 (\R/\Z)\times U\longrightarrow U,
 \qquad (a,w)\longmapsto\phi_{Ta}(w)
\]
is free.  Its quotient and quotient map will be denoted by
\[
 B=U/S^1,\qquad \rho:U\longrightarrow B.
\]
Thus $B$ is the closed smooth manifold of oriented, unmarked
geodesics, of dimension $2n-2$.

\subsection{The contact normalization and the two Gysin sequences}

The contact calculation is the volume method for periodic geodesic
flows associated with Weinstein \citep{Weinstein1974}; we include
it to fix the period, Euler sign and angular-area factors exactly.
The ensuing integral Gysin sequences, universal coefficient theorem
and Poincar\'e duality are used in their standard forms
\citep{Hatcher2002,MilnorStasheff1974}.
Yang's complex-projective volume calculation also uses the Gysin
sequences of the unit-tangent sphere bundle and the geodesic-circle
bundle \citep[pp.~381--382]{Yang1991}.  Here the degree-four
generator leads to a different middle-degree lattice.

Let $\lambda$ be the unit tangent contact form,
\[
 \lambda_w(\zeta)=\langle w,d\pi_w\zeta\rangle.
\]
Its Reeb flow is the unit-speed geodesic flow.  Consequently
$\theta=\lambda/T$ is a connection for the above $\R/\Z$ action:
it is invariant and integrates to one on each orbit.  Its curvature
is the pullback of a symplectic form $\omega_B$ on $B$.
We fix the Euler-class sign convention by
\[
 \rho^*\omega_B=d\theta,\qquad e=[\omega_B]\in H^2(B;\Z),
\]
and give $B$ its symplectic orientation.

\begin{proposition}[Contact volume]\label{geo:contact}
With this orientation and Euler class,
\begin{equation}\label{geo:contact-formula}
 N:=\int_B e^{4r-1}
   =\frac{(4r-1)!\,|S^{4r-1}|\,\Vol(M)}{T^{4r}}>0.
\end{equation}
For the standard quaternionic model at the same diameter, this
number is
\begin{equation}\label{geo:model-number}
 N_0=\frac{1}{2r+1}\binom{4r}{2r}.
\end{equation}
\end{proposition}

\begin{proof}
Integration over the circle fibers gives
\[
 \int_Be^d
 =\int_U\theta\wedge(d\theta)^d
 =T^{-(d+1)}\int_U\lambda\wedge(d\lambda)^d,
\]
where $U$ has the orientation induced by the positive circle
direction and the chosen orientation of $B$.
To compute the last form as a density, use the connection splitting
of $T_wU$ into horizontal and vertical directions.  Choose an
orthonormal frame $(w,e_1,\ldots,e_d)$ at $\pi(w)$.
The contact form reads off the horizontal $w$ coordinate, and
$d\lambda$ pairs each horizontal $e_i$ coordinate with its vertical
$e_i$ coordinate, up to one common sign convention.
In the $d$-fold exterior power each pairing occurs once in each of
$d!$ orders.  Hence
\[
 \big|\lambda\wedge(d\lambda)^d\big|=d!\,d\mu,
\]
where $d\mu$ is base Riemannian volume times unit tangent sphere
area.  The chosen orientation makes the contact integral positive,
and $\mu(U)=|S^d|\Vol(M)$.  Since $d+1=4r$, this proves
\eqref{geo:contact-formula}.

The model's polar density on $(0,D)$ is
$\sin^{4r-1}t\cos^3t$.  With $y=\sin^2t$,
\[
 \int_0^{\pi/2}\sin^{4r-1}t\cos^3t\,dt
 =\frac12\int_0^1 y^{2r-1}(1-y)\,dy
 =\frac{1}{4r(2r+1)}.
\]
Using $|S^{4r-1}|=2\pi^{2r}/(2r-1)!$ and $T=\pi$ in
\eqref{geo:contact-formula} gives
\[
 N_0=\frac{(4r-1)!}{r(2r+1)((2r-1)!)^2}
     =\frac{(4r)!}{(2r)!(2r+1)!},
\]
which is \eqref{geo:model-number}.
\end{proof}

Thus the desired upper volume bound is $N\le N_0$.  The integral
cohomology calculation below computes the determinant of a middle
pairing whose entries are top-degree products, including this
characteristic number.

\begin{lemma}[Integral Gysin information]\label{geo:gysin}
The manifolds $U$ and $B$ are simply connected.  There is an
integral class $v\in H^4(B;\Z)$ with $\rho^*v=z$.
For $2j<n$, the following is an integral basis:
\begin{equation}\label{geo:low-basis}
 H^{2j}(B;\Z):\qquad
 e^{j-2a}v^a,\qquad 0\le a\le\lfloor j/2\rfloor.
\end{equation}
All odd cohomology of $B$ vanishes, and its even cohomology is free
abelian.  In particular $H^2(B;\Z)=\Z e$ and
$H^4(B;\Z)=\Z e^2\oplus\Z v$.
The rational cohomology ring is generated by $e,v$, and
\begin{equation}\label{geo:total-betti}
 \dim_{\mathbb Q}H^*(B;\mathbb Q)=2r(r+1).
\end{equation}

Multiplication by $e$ at the middle fits into an exact sequence
\begin{equation}\label{geo:middle-sequence}
 0\longrightarrow H^{n-2}(B;\Z)
 \xrightarrow{\ e\smile\ }H^n(B;\Z)
 \longrightarrow\Z/(r+1)\longrightarrow0.
\end{equation}
Consequently the integral bilinear form
\begin{equation}\label{geo:middle-form}
 (a,b)\longmapsto\int_B eab,\qquad a,b\in H^{n-2}(B;\Z),
\end{equation}
has determinant of absolute value $r+1$ in any integral basis.
The map $e\smile:H^n(B;\Z)\to H^{n+2}(B;\Z)$ is an
isomorphism.
\end{lemma}

\begin{proof}
The sphere fiber of $\pi$ has dimension $n-1\ge7$.  Its homotopy
sequence and simple connectivity of $M$ give $\pi_1(U)=0$.
The homotopy sequence of the circle bundle then gives $\pi_1(B)=0$.

The Euler number of $TM$ is $\chi(M)=r+1$.  The chosen orientation
therefore gives
\[
 e(TM)=(r+1)u^r.
\]
The integral Gysin sequence of the oriented sphere bundle $\pi$
has the form
\[
 H^{q-n}(M)\xrightarrow{\ e(TM)\smile\ }H^q(M)
 \xrightarrow{\ \pi^*\ }H^q(U)
 \xrightarrow{\ \pi_!\ }H^{q-n+1}(M)
 \xrightarrow{\ e(TM)\smile\ }H^{q+1}(M).
\]
It gives the complete additive calculation
\begin{equation}\label{geo:unit-cohomology}
 H^q(U;\Z)=
 \begin{cases}
  \Z\{z^a\},&q=4a,\quad 0\le a<r,\\
  \Z/(r+1)\{z^r\},&q=n,\\
  \Z\{y_a\},&q=n-1+4a,\quad 1\le a\le r,\\
  0,&\text{otherwise}.
 \end{cases}
\end{equation}
Here the upper generators may be chosen with
$\pi_!y_a=u^a$ and $y_a=z^{a-1}y_1$; the latter assertion follows
from the projection formula and the fact that $\pi_!$ is an
isomorphism in these upper degrees.
At degree $n-1$, the Gysin map from $H^0(M)$ to $H^n(M)$ is
multiplication by the nonzero integer $r+1$, so
$H^{n-1}(U)=0$.  This also explains the sole torsion group in
\eqref{geo:unit-cohomology}.

Apply next the integral Gysin sequence of $\rho$.  The vanishing
of odd groups of $U$ below degree $n$ inductively gives
$H^{2j-1}(B)=0$ for $2j-1<n$.  For positive even $2j<n$ the
remaining terms give
\begin{equation}\label{geo:low-exact}
 0\longrightarrow H^{2j-2}(B)
 \xrightarrow{\ e\smile\ }H^{2j}(B)
 \xrightarrow{\ \rho^*\ }H^{2j}(U)
 \longrightarrow0.
\end{equation}
For $2j=2$ this says $H^2(B)=\Z e$.  For $2j=4$, the last
group is freely generated by $z$, so choose an integral lift $v$.
If $j$ is odd, the last group in \eqref{geo:low-exact} vanishes
and multiplication by $e$ gives the next basis.  If $j=2a$ is
even, the last group is generated by $z^a$, and the class $v^a$
lifts that generator.  Induction proves exactly the integral basis
\eqref{geo:low-basis}.

In all degrees at most $n-2$, the homology groups of $B$ are free
and have the same ranks as its cohomology groups.  Indeed the
universal coefficient theorem detects torsion in $H_i(B;\Z)$
as torsion in $H^{i+1}(B;\Z)$, and the latter groups are free
through degree $n-1$.  Poincar\'e duality on the oriented
$(2n-2)$-manifold $B$ now supplies all its cohomology in degrees
at least $n$.  It is free in even degrees and zero in odd degrees;
together with the low-degree calculation this proves the asserted
global statement.

The circle Gysin sequence at degree $n$ uses
$H^{n-1}(U)=H^{n-1}(B)=0$ and
$H^n(U)=\Z/(r+1)$, and is precisely
\eqref{geo:middle-sequence}.  The Poincar\'e pairing
\[
 H^{n-2}(B;\Z)\times H^n(B;\Z)\longrightarrow\Z
\]
is unimodular, since these two groups are free.  Composing one
factor with the injection of index $r+1$ in
\eqref{geo:middle-sequence} proves the determinant assertion
\eqref{geo:middle-form}.
Finally $H^{n+1}(U)=H^{n+2}(U)=0$ by
\eqref{geo:unit-cohomology}.  The next portion of the circle
Gysin sequence therefore makes multiplication by $e$ from degree
$n$ to degree $n+2$ an isomorphism.

In every even degree $q\ge n$, the group $H^q(U;\mathbb Q)$
vanishes.  The rational circle Gysin sequence therefore makes
$e\smile:H^{q-2}(B;\mathbb Q)\to H^q(B;\mathbb Q)$
surjective.  The low-degree bases now show that the full rational
ring is generated by $e,v$.  Below degree $n$, the basis ranks
sum to
\[
 \sum_{j=0}^{2r-1}\bigl(\lfloor j/2\rfloor+1\bigr)
 =2\sum_{a=0}^{r-1}(a+1)=r(r+1).
\]
Poincar\'e duality supplies another $r(r+1)$ dimensions above
degree $n-1$, and the middle odd group is zero.  This proves
\eqref{geo:total-betti}.
\end{proof}

\subsection{The first free-loop attachment}

The Gysin sequences determine the low-degree integral lattice and its middle
determinant, but not the relations among $e$ and $v$.  Those
relations will come from the first nonconstant critical manifold
of the free-loop energy.

Write $\Lambda M=H^1(\R/T\Z,M)$ for the Sobolev free loop
space, with rotation of the parameter as its circle action.
The continuous-loop and Sobolev-loop models have the same homotopy
type, compatibly with evaluation and rotation.  We use the energy
\[
 \mathcal E(c)=\frac12\int_0^T|\dot c(t)|^2\,dt.
\]
For a circle space $A$, write $A_{hS^1}=ES^1\times_{S^1}A$.
The constant-loop subspace has trivial action, hence its Borel
construction is $M\times BS^1$.

\begin{proposition}[The first attachment and its coefficient]
\label{geo:attachment}
There is an oriented rank-three vector bundle $E\to B$ and an
attaching map
\[
 f:S(E)\longrightarrow M.
\]
Let $p:S(E)\to B$ and put $x=f^*u\in H^4(S(E);\Z)$.
The following properties hold:
\begin{enumerate}[label=\textup{(\roman*)}]
\item The pullback $\rho^*E$ is the first negative bundle for the
      free-loop energy.  It is isomorphic as a real vector bundle
      to the rank-three bundle $K\to U$ with fiber
      \[
       K_w=\{a\in w^\perp:
          J_a(D)=0\},
       \quad J_a(0)=0,\quad J_a'(0)=a,
      \]
      where $J_a$ is the normal Jacobi field along $\gamma_w$.
\item Let $\widetilde p:S(\rho^*E)\to U$ and
      $\widehat\rho:S(\rho^*E)\to S(E)$ be the pullback square.
      Then
      \begin{equation}\label{geo:attachment-pullback}
       \widehat\rho^*x=\widetilde p^*z.
      \end{equation}
\item After choosing the orientation of $E$, one has
      \begin{equation}\label{geo:primitive-boundary}
       p_!x=e.
      \end{equation}
      With either orientation its coefficient is $\pm1$.
\item The actual negative bundle satisfies
      \begin{equation}\label{geo:negative-pontryagin}
       \rho^*p_1(E)=4z.
      \end{equation}
\end{enumerate}
\end{proposition}

The attaching map makes $x^{r+1}=f^*(u^{r+1})=0$.
The sphere integral $p_!x=e$ fixes the coefficient of the fiber
part of $x$, while $\rho^*p_1(E)=4z$ carries the incidence
coefficient to the actual negative bundle.  Separating $x$ into
base and fiber parts will turn its nilpotence into relations on
$B$.  The middle determinant will then fix the scale of the
top-degree pairing used to evaluate $N$.

Figure~\ref{fig:attachment-pullback} distinguishes the pullback
square from the evaluation homotopy used in part~\textup{(ii)}.
\begin{figure}[tbp]
\centering
\begin{tikzpicture}[>=Stealth,font=\small]
  \node (SU) at (-3,1.8) {$S(\rho^*E)$};
  \node (SB) at (1,1.8) {$S(E)$};
  \node (U) at (-3,-.2) {$U=SM$};
  \node (B) at (1,-.2) {$B=U/S^1$};
  \node (M) at (4,1.8) {$M$};
  \draw[->] (SU)--node[above] {$\widehat\rho$} (SB);
  \draw[->] (SB)--node[above] {$f$} (M);
  \draw[->] (SU)--node[left,align=right] {$\widetilde p$\\$S^2\text{ fibers}$} (U);
  \draw[->] (SB)--node[right,align=left] {$p$\\$S^2\text{ fibers}$} (B);
  \draw[->] (U)--node[below] {$\rho\quad(S^1\text{ fibers})$} (B);
  \node[fill=white] at (-1,.8) {pullback};
  \node[align=center] at (-3,-1.15) {marked unit-speed\\oriented geodesic};
  \node[align=center] at (1,-1.15) {oriented geodesic\\with no marked point};
  \node[align=center] at (.2,-2.25)
    {$f\circ\widehat\rho\simeq\pi\circ\widetilde p$
      \quad (evaluation homotopy)};
  \node[align=center] at (.2,-3.0)
    {$x=f^*u,\quad z=\pi^*u,\quad
       \widehat\rho^*x=\widetilde p^*z$};
\end{tikzpicture}

\caption{The sphere bundle of the descended rank-three bundle $E$
and its pullback to $U=SM$. The square is a pullback, with vertical fibers
$S^2$ and horizontal fibers $S^1$. The attaching map $f$ is not
point evaluation of an unmarked geodesic. Evaluation of the
attachment gives the displayed homotopy and hence the equality
of cohomology classes.}
\label{fig:attachment-pullback}
\end{figure}

\begin{proof}
We give the index, attachment, and integral-coefficient details.

\paragraph{Critical manifolds and indices.}
The positive critical levels are $j^2T/2$, $j\ge1$; the critical
manifold at each is a copy of $U$, parametrized by
$w\mapsto(t\mapsto\gamma_w(jt))$.  These are all the nonconstant
critical loops because every geodesic has prime length $T$.

For the $j$th iterate, reparametrize variations by arclength on
$[0,jT]$.  The kernel of the periodic Hessian consists of all
periodic Jacobi fields.  All normal Jacobi fields are $T$-periodic
by the differential of $\phi_T=\Id$, so their dimension is $2d$.
A tangential Jacobi field is an affine scalar multiple of the
velocity; the periodic boundary conditions force its scalar to be
constant.  The kernel therefore has dimension $2d+1=\dim U$.
Variations of the initial unit tangent vector, including the
time-shift direction, realize precisely these fields.
Thus the Hessian kernel is the
tangent space to the critical manifold, as required by the
Morse--Bott condition.

The periodic index equals the fixed-endpoint index.  To see this
without an index correction formula, evaluation at the marked
initial point maps the periodic kernel onto $T_{\pi(w)}M$:
arbitrary normal initial values are possible, and the constant
tangent field supplies the remaining direction.  Choose a linear
right inverse to this evaluation.  Subtracting the corresponding
kernel field from a periodic test field makes both endpoint
values zero.  The index form is unchanged, including its mixed
terms, because a kernel field pairs to zero with every periodic
test field by integration by parts.  This subtraction is injective
on every negative definite subspace.  Conversely, fixed-endpoint
test fields form a subspace of the periodic domain.  The two
indices are consequently equal.

Interior conjugate instants on $[0,jT]$ occur at the $j$ odd
multiples of $D$, with multiplicity three, and the $j-1$ positive
multiples of $T$, with multiplicity $d$.  The fixed-endpoint Morse
index theorem \citep[Theorem~15.1, p.~83]{Milnor1963} gives
\begin{equation}\label{geo:iterate-index}
 \lambda_j=3j+d(j-1).
\end{equation}
In particular $\lambda_1=3$, while every later index is at least
$d+6=n+5\ge13$. The index count is unchanged on completing the
piecewise smooth zero-endpoint fields in $H^1$: the index form
is continuous, and any finite negative Gram matrix remains
negative after sufficiently close smooth approximation.

\paragraph{Negative bundles and attachment.}
Let $\nu\to U$ be the negative spectral bundle of the periodic
$L^2$ Jacobi operator
\[
 \mathscr A_w V=-\nabla_t^2V-R(V,\dot\gamma_w)\dot\gamma_w.
\]
This self-adjoint operator has domain the periodic $H^2$ fields,
and its quadratic form is the index form on periodic $H^1$ fields.
Its negative eigenvalues form an isolated finite cluster. The
constant index and nullity, together with compactness of $U$,
make the corresponding spectral projections locally smooth in
the initial data. Elliptic regularity makes their eigenfields
smooth, so $\nu$ is a genuine rank-three bundle.
It represents the negative bundle for the Hilbert-space
Morse--Bott attachment as well. Indeed, projection from the
negative subspace of the bounded $H^1$-Riesz Hessian to $\nu_w$
is injective: a vector in its kernel would have nonnegative
index form in the $L^2$ spectral decomposition. Equal negative
ranks make this a bundle isomorphism.
Each fixed parameter rotation intertwines these operators and
projections with those at the rotated geodesic. The resulting
action on $\nu$ is continuous and covers the free smooth circle
action on $U$. Hence $\nu$ descends to a real bundle $E\to B$,
with $\nu\cong\rho^*E$. Since $B$ is simply connected, $E$ is
orientable.

We use the Morse--Bott disk-attachment theorem for the energy on
the $H^1$ loop space of a closed Riemannian manifold.  The energy
satisfies the Palais--Smale condition.  Its Hessian is nondegenerate
normal to each compact critical manifold just described, and has
finite negative index. The bounded $H^1$ Hessian is Fredholm;
its identified kernel therefore gives an invertible restriction
normal to the critical manifold, as in the formulation cited
below. Thus sublevels bracketing one critical
value differ, up to homotopy, by the disk and sphere bundles of
its negative bundle.  This construction also respects parameter
rotation.  Indeed each fixed rotation is a smooth isometry for
the natural loop-space metric, preserves energy, commutes with
the gradient deformation, and preserves the Hessian splitting.
The local normal form and the contraction of its positive
directions may be constructed from these invariant metrics and
splittings.  Joint differentiability of the circle action on the
whole $H^1$ manifold is not needed.  For this equivariant
negative-bundle formulation see \citet[Section~7, pp.~2190--2192]
{OttosenBokstedt2007}, following the Hilbert-space attachment of
\citet[Theorem~2.4.10]{Klingenberg1978}.
The associated Borel vector bundle and Thom quotient are described
in \citet[Lemma~5.1, pp.~2185--2186]{OttosenBokstedt2007}.

Here is the attaching map and its evaluation property explicitly.
Start a sufficiently small negative disk at every first-level
geodesic by pointwise exponential of the negative fields.  Its
boundary lies strictly below the first positive critical value,
uniformly over $U$.  That lower sublevel deforms equivariantly to
the constants.  This follows from the noncritical energy
deformation between the boundary level and a small positive
level, followed by contraction of small loops to constants.  For
the latter, embed $M$ in Euclidean space, average a sufficiently
short loop, and project the straight contraction to a tubular
neighborhood of $M$.  Small energy gives uniformly small length,
so this is defined and continuous, preserves the constants, and
commutes with rotation.  The projection has bounded derivative,
so choosing the initial small level sufficiently small keeps this
contraction below the first critical value.

Insert the boundary deformation in an outer collar of the small
negative disk bundle.  This produces an equivariant map of pairs
\[
 \mathcal F_1:(D(\nu),S(\nu))\longrightarrow(\Lambda M,M)
\]
whose zero section is the first critical manifold and which
represents its Morse--Bott attachment.  On the boundary it takes
values in the constants.  Because the action on the constants
is trivial, this boundary map descends to
$f:S(E)\to M$.  Evaluation at time zero composed with
$\mathcal F_1$ is homotopic, by radial contraction of the disk
fibers, to $\pi$ composed with their bundle projection.  Restrict
this homotopy to $S(\nu)\cong S(\rho^*E)$.  Its boundary value
is $f\circ\widehat\rho$, which proves
\eqref{geo:attachment-pullback}.

Choose regular sublevels $\Lambda_-$ and $\Lambda_+$ immediately
below and above the first positive critical value.  The preceding
attachment gives, over $\Z$ or any field $\mathbb F$,
\begin{align}
 H^q(\Lambda_+,M;\mathbb F)
     &\cong H^{q-3}(U;\mathbb F),
       \label{geo:ordinary-relative}\\
 H^q((\Lambda_+)_{hS^1},M\times BS^1;\mathbb F)
     &\cong H^{q-3}(B;\mathbb F).
       \label{geo:borel-relative}
\end{align}
These are Thom isomorphisms; in the second one the free first
critical action identifies the Borel negative bundle with $E$
up to homotopy.  One can take finite-dimensional free
approximations to $ES^1$ throughout, apply the invariant
attachment, and then pass to the limit.
For every later critical level, the Borel negative bundle still
exists over the Borel critical manifold even though the action
has finite stabilizers.  Its relative cohomology vanishes below
its rank $\lambda_j$, with any necessary orientation local
system.  The same statement holds without the Borel construction.
Equation~\eqref{geo:iterate-index} therefore shows that subsequent
attachments do not change cohomology in degrees at most five.
The exhaustion by sublevels causes no extra term in these degrees:
the restriction systems in the relevant degrees are eventually
constant, so their inverse-limit obstruction vanishes.

We have constructed the attachment and isolated the degrees in
which no later attachment can alter it. What remains is to fix
its coefficient integrally and identify its actual negative bundle.

\paragraph{The coefficient is a unit over every prime field.}
Since $H^2(B;\Z)=\Z e$, write $p_!x=c e$ with $c\in\Z$.
Under \eqref{geo:borel-relative}, the connecting homomorphism
\[
 H^4(M\times BS^1;\mathbb F)
 \longrightarrow
 H^5((\Lambda_+)_{hS^1},M\times BS^1;\mathbb F)
 \cong H^2(B;\mathbb F)
\]
sends the class $u$ from $M$ to $p_!x$, up to the Thom
orientation sign.  Indeed the connecting map first restricts
to the attaching sphere, where that class is $x$, and then
uses the boundary map for $(D(E),S(E))$; the inverse Thom
isomorphism identifies this last map with integration over the
oriented sphere fiber.

Suppose that $c$ vanished modulo a prime $\ell$, and put
$\mathbb F=\Z/\ell$.  Exactness would extend the pure class
$u\in H^4(M\times BS^1;\mathbb F)$ over
$(\Lambda_+)_{hS^1}$ and hence over $(\Lambda M)_{hS^1}$,
by the higher-index gap.  After forgetting equivariance, its
restriction would be the evaluation class
\[
 a=\operatorname{ev}_0^*u\in H^4(\Lambda M;\mathbb F).
\]
To justify uniqueness here, \eqref{geo:ordinary-relative} gives
$H^4(\Lambda_+,M;\mathbb F)=H^1(U;\mathbb F)=0$.
Higher attachments have no effect in this range, so restriction
$H^4(\Lambda M;\mathbb F)\to H^4(M;\mathbb F)$ is injective.
The evaluation class is already one extension of $u$, and thus
is the only one.

An ordinary class restricted from the Borel construction has
zero action slant.  For completeness, if $Y$ is any circle space,
fix $e_0\in ES^1$ and let $i(y)=[e_0,y]\in Y_{hS^1}$.
For the action map $A:S^1\times Y\to Y$, the two maps
$i\circ A$ and $i\circ\operatorname{pr}_Y$ are homotopic:
\[
 [e_0,gy]=[e_0g,y],
\]
and the loop $g\mapsto e_0g$ contracts in $ES^1$.
Consequently the pullback of a Borel class under $i\circ A$
has no component in
$H^1(S^1;\mathbb F)\otimes H^{*-1}(Y;\mathbb F)$.
Equivalently its slant with $[S^1]$ is zero.  This is also the
usual action description of the Borel $d_2$ differential.

The evaluation class $a$ has a nonzero action slant over every
prime field.  The integral cohomology ring of $M$ and the
universal coefficient theorem give $H_2(M;\Z)=H_3(M;\Z)=0$.
Simple connectivity and successive applications of the Hurewicz
theorem \citep[Theorem~4.32]{Hatcher2002} make $M$
three-connected; Hurewicz then identifies
$\pi_4(M)$ with $H_4(M;\Z)=\Z$.
Choose a based map $g:S^4\to M$ with
$\langle g^*u,[S^4]\rangle=1$, and let
$\widehat g:S^3\to\Omega M\subset\Lambda M$ be its adjoint.
Evaluation of the rotated based loops on $S^1\times S^3$
factors, up to the orientation convention, through the smash
quotient
\[
 S^1\times S^3\longrightarrow S^1\wedge S^3=S^4
 \xrightarrow{\ g\ }M.
\]
The quotient takes the cross-product fundamental class to the
fundamental class of $S^4$, up to sign.  Therefore
\[
 \left\langle\widehat g^*
       \big(A^*\operatorname{ev}_0^*u/[S^1]\big),[S^3]\right\rangle
 =\pm1
\]
integrally, and remains nonzero after reduction modulo $\ell$.
This contradicts the zero slant required of the hypothesized
Borel extension.  No prime divides $c$, so $c=\pm1$.
Reverse the orientation of $E$ if necessary to obtain
\eqref{geo:primitive-boundary}.

\paragraph{The negative bundle is the cut-fiber tangent bundle.}
For $a\in K_w$, truncate $J_a$ to the first half-period and
extend it by zero on the second half:
\[
 \widehat J_a(t)=
 \begin{cases}
 J_a(t),&0\le t\le D,\\
 0,&D\le t\le T.
 \end{cases}
\]
The endpoint values agree, so this is an $H^1$ periodic test
field.  Integration by parts on $[0,D]$, with the vanishing
endpoint values and the Jacobi equation, gives zero index form
on this entire three-dimensional space.  Project these fields
to the negative spectral space $\nu_w$.  If a projected vector
were zero, its zero index form and its membership in the
nonnegative spectral subspace would force it into the Hessian
kernel.  It would then be a smooth periodic Jacobi field.
But a nonzero $J_a$ has nonzero derivative at $D$, since both
zero value and zero derivative would make it identically zero.
Its truncation has a derivative jump at $D$ and cannot be a
kernel field.  The projection is thus injective, and equal ranks
make it an isomorphism $K_w\to\nu_w$.
The truncation depends continuously in the $H^1$ norm on $w$
and $a$, and the spectral projections vary continuously as
above.  We obtain an isomorphism of real bundles
$K\cong\nu\cong\rho^*E$.
The truncations are used in the $H^1$ form domain, not in the
$H^2$ operator domain; their derivative jump is the reason
that a nonzero truncation cannot lie in the smooth kernel.

Finally define the full cut map
\[
 \mathfrak a:U\longrightarrow F,
 \qquad \mathfrak a(w)=(\pi(w),\exp_{\pi(w)}(Dw)).
\]
By the actual bundle identification in
Proposition~\ref{inc:normal}, its pullback of $\mathcal V$
is the initial-direction four-plane.  The marked unit direction
$w$ gives a trivial line in this plane; its orthogonal complement
is exactly $K_w$, the tangent space of the great-three-sphere
fiber.  Hence
\[
 \mathfrak a^*\mathcal V\cong\R\oplus K
 \cong\R\oplus\rho^*E.
\]
The two point maps from $U$ to $M$ in $\mathfrak a$ are
$\pi$ and $\pi\circ\phi_D$.  The geodesic flow provides a
homotopy between them, so
$\mathfrak a^*u_1=\mathfrak a^*u_2=z$.
Taking first Pontryagin classes and using
\eqref{geo:incidence-input} proves
\eqref{geo:negative-pontryagin}.
\end{proof}

\subsection{Half-integral lifts and the middle pairing}

On each sphere fiber, half the Euler class of its tangent bundle
integrates to one.  This gives a normalized fiber class with which
to separate $x$ into base and fiber parts.  The base class may be
only half-integral; we will track its displacement from the integral
lift $v$ so that the middle determinant is preserved.

\begin{proposition}[Arithmetic parameters]\label{geo:parameters}
Let $p:S(E)\to B$ and $x$ be as in
Proposition~\ref{geo:attachment}, and let
\[
 \xi=\frac12e(T_{\mathrm{vert}}S(E))\in H^2(S(E);\mathbb Q),
 \qquad X=e^2\in H^4(B;\Z).
\]
There is a class $Y\in H^4(B;\mathbb Q)$ such that, with
pullbacks by $p$ suppressed,
\begin{equation}\label{geo:sphere-parameters}
 x=Y+e\xi,\qquad
 \xi^2=p_1(E)/4,\qquad p_!\xi=1,\qquad
 \rho^*Y=z,\qquad 2Y\in H^4(B;\Z).
\end{equation}
For an integer $b$,
\begin{equation}\label{geo:Y-lift}
 Y=v+\frac b2e^2.
\end{equation}
There are rational numbers $\alpha,\beta$ with
\begin{equation}\label{geo:L-parameters}
 L:=e^2p_1(E)/4=\alpha X^2+\beta XY,\qquad
 4\alpha\in\Z,\qquad |\beta|=1.
\end{equation}
Here $\beta=1$ for \eqref{geo:incidence-input}; retaining $\beta$
also allows the opposite generator convention, for which
$\beta=-1$.
Moreover
\begin{equation}\label{geo:two-nilpotents}
 (Y+e\xi)^{r+1}=(Y-e\xi)^{r+1}=0
 \quad\text{in }H^*(S(E);\mathbb Q).
\end{equation}

Count $X,Y$ as polynomial degree one, corresponding to
cohomological degree four.  For homogeneous polynomials of
degree $2r-1$, set
\[
 \ell(P)=\int_B eP(X,Y).
\]
Then the middle moment matrix satisfies
\begin{equation}\label{geo:moment-determinant}
 \left|\det\left[
 \ell\big(X^{2r-1-i-j}Y^{i+j}\big)
 \right]_{0\le i,j<r}\right|=r+1,
 \qquad \ell(X^{2r-1})=N.
\end{equation}
Pullback $p^*:H^*(B;\mathbb Q)\to H^*(S(E);\mathbb Q)$
is injective, and multiplication by $e$ on $H^n(B;\mathbb Q)$
is injective.
\end{proposition}

\begin{proof}
Give each unit sphere fiber the boundary orientation from $E$.
The marked unit vector in $p^*E$ spans a trivial line, with
orthogonal complement the oriented vertical tangent two-plane:
\[
 p^*E\cong\R\oplus T_{\mathrm{vert}}S(E).
\]
For an oriented real two-plane bundle, the first Pontryagin
class is the square of its Euler class.  Thus
$\xi^2=p^*p_1(E)/4$.  On a sphere fiber the Euler class of
its tangent bundle integrates to two, giving $p_!\xi=1$.

The Euler class of the oriented rank-three bundle $E$ lies in
$H^3(B;\Z)=0$, so it vanishes.  The sphere-bundle Gysin
sequence then gives, in every degree, a short exact sequence
whose first map is $p^*$ and whose last map is $p_!$.
In particular $p^*$ is injective.  Over $\mathbb Q$, the
classes $1,\xi$ restrict to a basis of each sphere's cohomology,
so Leray--Hirsch gives the module decomposition
\[
 H^*(S(E);\mathbb Q)
 =p^*H^*(B;\mathbb Q)
  \oplus\xi\,p^*H^{*-2}(B;\mathbb Q).
\]
Since $p_!x=e$, there is a unique rational class $Y$ with
$x=p^*Y+p^*e\,\xi$.
The integral class
\[
 2x-p^*e\,e(T_{\mathrm{vert}}S(E))
\]
has zero sphere integral.  Integral Gysin exactness expresses
it as $p^*Y_2$ for $Y_2\in H^4(B;\Z)$.  Rational injectivity
of $p^*$ gives $Y_2=2Y$.  This proves half-integrality without
assuming that $Y$ itself is integral.

Pull back to the square over $U$ in
Proposition~\ref{geo:attachment}.  The class $\rho^*e$ vanishes,
as it does for the Euler class pulled to its own circle bundle.
Equation~\eqref{geo:attachment-pullback} now gives
\[
 \widetilde p^*(\rho^*Y)=\widehat\rho^*x
 =\widetilde p^*z.
\]
The pullback sphere bundle also has zero Euler class; its
rational Gysin sequence makes $\widetilde p^*$ injective.
It follows that $\rho^*Y=z$.
Write $2Y=a v+b e^2$ in the integral basis supplied by
Lemma~\ref{geo:gysin}.  Since $\rho^*v=z$ is a primitive
generator of $H^4(U;\Z)$, pulling back gives $a=2$.
This proves \eqref{geo:Y-lift}.

Similarly, integrality of $p_1(E)$ and
\eqref{geo:negative-pontryagin} give
\[
 p_1(E)=A e^2+4v,\qquad A\in\Z.
\]
If the opposite sign of the incidence coefficient is retained,
this reads $p_1(E)=A e^2+4\beta v$, with $\beta=\pm1$.
In either convention substitution of
$v=Y-(b/2)e^2$ gives
\[
 L=\frac{A-2\beta b}{4}X^2+\beta XY.
\]
Thus $4\alpha=A-2\beta b$ is integral, proving
\eqref{geo:L-parameters}.

The first equality in \eqref{geo:two-nilpotents} is
$x^{r+1}=f^*(u^{r+1})=0$.  Fiber antipodal reflection covers
the identity of $B$ and reverses the orientation of the
two-dimensional sphere fibers.  Its pullback fixes base classes
and sends the vertical Euler class, hence $\xi$, to its negative.
Pulling back the first equality proves the second one.

It remains to identify the moment determinant.  An integral
basis of $H^{n-2}(B;\Z)$ from
\eqref{geo:low-basis} is
\[
 b_i=eX^{r-1-i}v^i,\qquad 0\le i<r.
\]
The matrix of \eqref{geo:middle-form} in this basis has entries
\[
 \int_Be b_i b_j
 =\int_B eX^{2r-1-i-j}v^{i+j}.
\]
Replace $v$ in these basis vectors by
$Y=v+(b/2)X$.  The change of basis over $\mathbb Q$ is triangular
with every diagonal entry equal to one.  Hence its determinant
is one and it does not change the determinant of the bilinear
form.  Lemma~\ref{geo:gysin} proves the first assertion in
\eqref{geo:moment-determinant}.  The second follows from
$eX^{2r-1}=e^{4r-1}$ and the definition of $N$.
The final injectivity statements were proved above for $p^*$
and in Lemma~\ref{geo:gysin} for multiplication by $e$.
\end{proof}
